% !TeX program = lualatex
% !TeX root = predavanje11.tex
% ============================================================
% Zajednički preamble za sva predavanja RWA kolegija
% Uključiti s: % ============================================================
% Zajednički preamble za sva predavanja RWA kolegija
% Uključiti s: % ============================================================
% Zajednički preamble za sva predavanja RWA kolegija
% Uključiti s: \input{../common/preamble.tex}
% ============================================================

\documentclass[aspectratio=149,11pt]{beamer}
\usetheme{CambridgeUS}

\usepackage{fontspec}              % radi samo sa XeLaTeX/LuaLaTeX
\setsansfont{Latin Modern Sans}    % sans-serif font
\setmonofont{Latin Modern Mono}    % monospace font za kod
\usepackage[croatian]{babel}
\usepackage{csquotes}              % \enquote{} → hrvatski navodnici „..."
\usepackage{lmodern}
\usepackage{booktabs}
\usepackage{listings}
\usepackage{xcolor}
\usepackage{hyperref}

% ---------- Autor i logo ----------
\author{mag.ing. Josip Torić (SIT UNIZD)}
\titlegraphic{\includegraphics[height=2.5cm]{../common/SIT_logo_tekst1.pdf}}

% ---------- Listings ----------
\definecolor{codebg}{RGB}{248,248,248}
\definecolor{codeframeclr}{RGB}{220,220,220}
\lstdefinelanguage{bash}{
  morekeywords={sudo,apt,brew,pip,pipx,uv,winget,python,git,flask,source,deactivate,code,gunicorn,docker,compose,heroku,render,railway,node,npm,pnpm,yarn,uvicorn,alembic,pytest,ruff,black,precommit},
  sensitive=false,
  morecomment=[l]{\#}
}
\lstdefinelanguage{python}{
  morekeywords={from,import,def,class,return,if,elif,else,for,in,while,try,except,raise,with,as,not,and,or,is,lambda,yield,True,False,None,await,async},
  sensitive=true,
  morecomment=[l]{\#},
  morestring=[b]',
  morestring=[b]"
}
\lstdefinelanguage{sql}{
  morekeywords={SELECT,FROM,WHERE,INSERT,INTO,VALUES,UPDATE,SET,DELETE,CREATE,TABLE,PRIMARY,KEY,FOREIGN,REFERENCES,UNIQUE,NOT,NULL,INDEX,ON,ALTER,DROP,CONSTRAINT},
  sensitive=false,
  morecomment=[l]{--}
}
\lstdefinelanguage{json}{
  morecomment=[l]{//},
  morestring=[b]",
  showstringspaces=false
}
\lstdefinelanguage{typescript}{
  morekeywords={import,export,from,const,let,var,function,return,if,else,for,of,in,while,switch,case,break,async,await,try,catch,finally,throw,new,class,this,extends,implements,interface,type,enum,readonly,public,private,protected,static,get,set,true,false,null,undefined,void,string,number,boolean,any,unknown,never,typeof,keyof,as},
  sensitive=true,
  morecomment=[l]{//},
  morecomment=[s]{/*}{*/},
  morestring=[b]',
  morestring=[b]",
  morestring=[b]`
}
\lstalias{javascript}{typescript}
\lstset{
  backgroundcolor=\color{codebg},
  basicstyle=\ttfamily\footnotesize,
  frame=single,
  rulecolor=\color{codeframeclr},
  frameround=tttt,
  breaklines=true,
  tabsize=2,
  showstringspaces=false,
  upquote=true,
  columns=fullflexible
}

% ---------- Helper macros ----------
\newcommand{\codefile}[1]{\texttt{#1}}
\newcommand{\kbd}[1]{\texttt{#1}}

% ---------- Code frame environment ----------
% Automatski dodaje fragile i shrink=5, koristi se umjesto \begin{frame}
% Primjer: \begin{codeframe}{Naslov} ... \end{codeframe}
% Primjer s opcijama: \begin{codeframe}[shrink=15]{Naslov} ... \end{codeframe}
\newenvironment{codeframe}[2][]
  {\begin{frame}[fragile,shrink=5,#1]{#2}}
  {\end{frame}}

% ============================================================

\documentclass[aspectratio=149,11pt]{beamer}
\usetheme{CambridgeUS}

\usepackage{fontspec}              % radi samo sa XeLaTeX/LuaLaTeX
\setsansfont{Latin Modern Sans}    % sans-serif font
\setmonofont{Latin Modern Mono}    % monospace font za kod
\usepackage[croatian]{babel}
\usepackage{csquotes}              % \enquote{} → hrvatski navodnici „..."
\usepackage{lmodern}
\usepackage{booktabs}
\usepackage{listings}
\usepackage{xcolor}
\usepackage{hyperref}

% ---------- Autor i logo ----------
\author{mag.ing. Josip Torić (SIT UNIZD)}
\titlegraphic{\includegraphics[height=2.5cm]{../common/SIT_logo_tekst1.pdf}}

% ---------- Listings ----------
\definecolor{codebg}{RGB}{248,248,248}
\definecolor{codeframeclr}{RGB}{220,220,220}
\lstdefinelanguage{bash}{
  morekeywords={sudo,apt,brew,pip,pipx,uv,winget,python,git,flask,source,deactivate,code,gunicorn,docker,compose,heroku,render,railway,node,npm,pnpm,yarn,uvicorn,alembic,pytest,ruff,black,precommit},
  sensitive=false,
  morecomment=[l]{\#}
}
\lstdefinelanguage{python}{
  morekeywords={from,import,def,class,return,if,elif,else,for,in,while,try,except,raise,with,as,not,and,or,is,lambda,yield,True,False,None,await,async},
  sensitive=true,
  morecomment=[l]{\#},
  morestring=[b]',
  morestring=[b]"
}
\lstdefinelanguage{sql}{
  morekeywords={SELECT,FROM,WHERE,INSERT,INTO,VALUES,UPDATE,SET,DELETE,CREATE,TABLE,PRIMARY,KEY,FOREIGN,REFERENCES,UNIQUE,NOT,NULL,INDEX,ON,ALTER,DROP,CONSTRAINT},
  sensitive=false,
  morecomment=[l]{--}
}
\lstdefinelanguage{json}{
  morecomment=[l]{//},
  morestring=[b]",
  showstringspaces=false
}
\lstdefinelanguage{typescript}{
  morekeywords={import,export,from,const,let,var,function,return,if,else,for,of,in,while,switch,case,break,async,await,try,catch,finally,throw,new,class,this,extends,implements,interface,type,enum,readonly,public,private,protected,static,get,set,true,false,null,undefined,void,string,number,boolean,any,unknown,never,typeof,keyof,as},
  sensitive=true,
  morecomment=[l]{//},
  morecomment=[s]{/*}{*/},
  morestring=[b]',
  morestring=[b]",
  morestring=[b]`
}
\lstalias{javascript}{typescript}
\lstset{
  backgroundcolor=\color{codebg},
  basicstyle=\ttfamily\footnotesize,
  frame=single,
  rulecolor=\color{codeframeclr},
  frameround=tttt,
  breaklines=true,
  tabsize=2,
  showstringspaces=false,
  upquote=true,
  columns=fullflexible
}

% ---------- Helper macros ----------
\newcommand{\codefile}[1]{\texttt{#1}}
\newcommand{\kbd}[1]{\texttt{#1}}

% ---------- Code frame environment ----------
% Automatski dodaje fragile i shrink=5, koristi se umjesto \begin{frame}
% Primjer: \begin{codeframe}{Naslov} ... \end{codeframe}
% Primjer s opcijama: \begin{codeframe}[shrink=15]{Naslov} ... \end{codeframe}
\newenvironment{codeframe}[2][]
  {\begin{frame}[fragile,shrink=5,#1]{#2}}
  {\end{frame}}

% ============================================================

\documentclass[aspectratio=149,11pt]{beamer}
\usetheme{CambridgeUS}

\usepackage{fontspec}              % radi samo sa XeLaTeX/LuaLaTeX
\setsansfont{Latin Modern Sans}    % sans-serif font
\setmonofont{Latin Modern Mono}    % monospace font za kod
\usepackage[croatian]{babel}
\usepackage{csquotes}              % \enquote{} → hrvatski navodnici „..."
\usepackage{lmodern}
\usepackage{booktabs}
\usepackage{listings}
\usepackage{xcolor}
\usepackage{hyperref}

% ---------- Autor i logo ----------
\author{mag.ing. Josip Torić (SIT UNIZD)}
\titlegraphic{\includegraphics[height=2.5cm]{../common/SIT_logo_tekst1.pdf}}

% ---------- Listings ----------
\definecolor{codebg}{RGB}{248,248,248}
\definecolor{codeframeclr}{RGB}{220,220,220}
\lstdefinelanguage{bash}{
  morekeywords={sudo,apt,brew,pip,pipx,uv,winget,python,git,flask,source,deactivate,code,gunicorn,docker,compose,heroku,render,railway,node,npm,pnpm,yarn,uvicorn,alembic,pytest,ruff,black,precommit},
  sensitive=false,
  morecomment=[l]{\#}
}
\lstdefinelanguage{python}{
  morekeywords={from,import,def,class,return,if,elif,else,for,in,while,try,except,raise,with,as,not,and,or,is,lambda,yield,True,False,None,await,async},
  sensitive=true,
  morecomment=[l]{\#},
  morestring=[b]',
  morestring=[b]"
}
\lstdefinelanguage{sql}{
  morekeywords={SELECT,FROM,WHERE,INSERT,INTO,VALUES,UPDATE,SET,DELETE,CREATE,TABLE,PRIMARY,KEY,FOREIGN,REFERENCES,UNIQUE,NOT,NULL,INDEX,ON,ALTER,DROP,CONSTRAINT},
  sensitive=false,
  morecomment=[l]{--}
}
\lstdefinelanguage{json}{
  morecomment=[l]{//},
  morestring=[b]",
  showstringspaces=false
}
\lstdefinelanguage{typescript}{
  morekeywords={import,export,from,const,let,var,function,return,if,else,for,of,in,while,switch,case,break,async,await,try,catch,finally,throw,new,class,this,extends,implements,interface,type,enum,readonly,public,private,protected,static,get,set,true,false,null,undefined,void,string,number,boolean,any,unknown,never,typeof,keyof,as},
  sensitive=true,
  morecomment=[l]{//},
  morecomment=[s]{/*}{*/},
  morestring=[b]',
  morestring=[b]",
  morestring=[b]`
}
\lstalias{javascript}{typescript}
\lstset{
  backgroundcolor=\color{codebg},
  basicstyle=\ttfamily\footnotesize,
  frame=single,
  rulecolor=\color{codeframeclr},
  frameround=tttt,
  breaklines=true,
  tabsize=2,
  showstringspaces=false,
  upquote=true,
  columns=fullflexible
}

% ---------- Helper macros ----------
\newcommand{\codefile}[1]{\texttt{#1}}
\newcommand{\kbd}[1]{\texttt{#1}}

% ---------- Code frame environment ----------
% Automatski dodaje fragile i shrink=5, koristi se umjesto \begin{frame}
% Primjer: \begin{codeframe}{Naslov} ... \end{codeframe}
% Primjer s opcijama: \begin{codeframe}[shrink=15]{Naslov} ... \end{codeframe}
\newenvironment{codeframe}[2][]
  {\begin{frame}[fragile,shrink=5,#1]{#2}}
  {\end{frame}}


\title[Programiranje za web]{Programiranje za web - Predavanje 11 \\ Testiranje full-stack aplikacije}
\date{21.12.2026.}

\begin{document}

\frame{\titlepage}

% ------------------------------------------------------------
\begin{frame}{Plan predavanja}
\begin{enumerate}
  \item Zašto testiramo i piramida testiranja
  \item pytest: fixtures, conftest i izolacija testova
  \item Testiranje vremenskih pravila (freezegun)
  \item Coverage: koliko je dovoljno
  \item Vitest: unit testovi za storeove i utile
  \item Testiranje komponenti (@vue/test-utils)
  \item E2E testovi s Playwrightom
  \item Testovi u CI-ju i commitovi
\end{enumerate}

\vspace{0.6em}
Na kraju predavanja: backend i frontend imaju testove koji čuvaju poslovna pravila, a vi znate koji tip testa gdje pripada.
\end{frame}

% ------------------------------------------------------------
\begin{frame}{Ishodi učenja}
Nakon ovog predavanja student može:
\begin{itemize}
  \item objasniti piramidu testiranja i odabrati pravi tip testa za problem
  \item napisati pytest fixture i zamijeniti dependency s \texttt{dependency\_overrides}
  \item izolirati testove transakcijom s rollbackom
  \item testirati vremenski ovisna pravila s freezegunom
  \item napisati Vitest unit test za Pinia store i utility funkciju
  \item napisati component test s \texttt{@vue/test-utils}
  \item napisati Playwright E2E test za login i pravilo faze
\end{itemize}
\end{frame}

% ------------------------------------------------------------
\begin{frame}{Kontekst: gdje smo u projektu?}
\begin{itemize}
  \item Backend (P1-P6): API s pravilima, testove smo pisali usput uz svaku temu
  \item Frontend (P7-P10): SPA s auth flowom, formama i admin ekranima - bez ijednog testa
  \item Danas: sustavno testiranje \textbf{cijelog stacka} prije deploya
\end{itemize}

\vspace{0.6em}
Sljedeći korak (Predavanje 12): deploy. Aplikaciju koja ide u produkciju bez testova ne možete mijenjati bez straha - svaka izmjena može nešto potiho razbiti.
\end{frame}

% ------------------------------------------------------------
\begin{frame}{}
\begin{center}
\Large Test nije dokaz da kod radi. \\[0.5em]
Test je dokaz da \textbf{pravilo vrijedi} - i da vrijedi i sutra, nakon što nešto promijenite.
\end{center}
\end{frame}

% ------------------------------------------------------------
\begin{frame}{Piramida testiranja}
\begin{tabular}{@{}llllp{3.2cm}@{}}
  \toprule
  \textbf{Razina} & \textbf{Broj} & \textbf{Brzina} & \textbf{Alat} & \textbf{Što testira} \\
  \midrule
  E2E & malo & minute & Playwright & Cijeli flow u browseru \\
  Integration & srednje & sekunde & pytest + httpx & Endpoint kroz sve slojeve \\
  Unit & puno & milisekunde & pytest, Vitest & Izolirana funkcija, store \\
  \bottomrule
\end{tabular}

\vspace{0.6em}
Što je test niže u piramidi, to je brži i precizniji - kad padne, odmah znate gdje je problem. E2E test hvata ono što unit ne može: integraciju svih dijelova. Zato ih imamo malo, ali na ključnim flowovima.

\vspace{0.4em}
Anti-pattern je obrnuta piramida: sve kroz spori E2E, ništa kroz unit. Takvi testovi traju predugo i padaju iz nejasnih razloga.
\end{frame}

% ------------------------------------------------------------
\begin{frame}{Što smo dosad testirali (rekapitulacija)}
\begin{tabular}{@{}llp{5.5cm}@{}}
  \toprule
  \textbf{Predavanje} & \textbf{Testovi} & \textbf{Što čuvaju} \\
  \midrule
  P3 & Login, refresh, istekli token & Autentikacija radi i ne otkriva korisnike \\
  P4 & Role i cross-club pristup & Klub ne vidi tuđe resurse \\
  P5 & CRUD, validacija, duplikati & Constrainti i 409/422 odgovori \\
  P6 & Faze s freezegunom & Rokovi zabranjuju akcije \\
  \bottomrule
\end{tabular}

\vspace{0.6em}
Te testove ste pisali uz gradivo. Danas ih slažemo u strukturu koja se lako održava: zajedničke fixture, izolacija, coverage - i dodajemo frontend koji dosad nije imao ništa.
\end{frame}

% ------------------------------------------------------------
\begin{frame}{Anatomija dobrog testa: Arrange - Act - Assert}
Svaki test ima tri dijela:

\vspace{0.4em}
\begin{tabular}{@{}lp{8.5cm}@{}}
  \toprule
  \textbf{Dio} & \textbf{Što radi} \\
  \midrule
  Arrange & Pripremi stanje: korisnik u bazi, token, natjecanje s rokom \\
  Act & Izvrši jednu akciju: pošalji request, pozovi funkciju \\
  Assert & Provjeri ishod: status kod, sadržaj odgovora, stanje u bazi \\
  \bottomrule
\end{tabular}

\vspace{0.6em}
Ime testa opisuje pravilo, ne implementaciju: \texttt{test\_klub\_ne\_moze\_prijaviti\_nakon\_prelim\_roka} govori što je pravilo. Kad taj test padne, iz imena znate koje se pravilo prekršilo.

\vspace{0.4em}
Jedan test = jedno pravilo. Test koji provjerava pet stvari padne, a vi ne znate koja od pet je kriva.
\end{frame}

% ------------------------------------------------------------
\begin{frame}[fragile,shrink=5]{pytest fixtures: gradivni blokovi testova}
Fixture je funkcija koja priprema ono što test treba - pytest je poziva automatski po imenu parametra:

\begin{lstlisting}[language=python]
import pytest

@pytest.fixture
def competition_open():
    return Competition(
        name="Otvoreno prvenstvo",
        prelim_deadline=datetime(2027, 3, 1, tzinfo=timezone.utc),
        final_deadline=datetime(2027, 3, 15, tzinfo=timezone.utc),
    )

def test_faza_prije_prelim_roka(competition_open):
    with freeze_time("2027-02-01"):
        assert get_competition_phase(competition_open) == Phase.OPEN
\end{lstlisting}

\vspace{0.4em}
Fixture se definira jednom, koristi u koliko god testova treba. Kad se model promijeni, mijenjate jednu fixture umjesto dvadeset testova.
\end{frame}

% ------------------------------------------------------------
\begin{frame}[fragile,shrink=10]{conftest.py: dijeljene fixture za cijeli projekt}
\begin{lstlisting}[language=python]
# tests/conftest.py
import pytest
from httpx import AsyncClient, ASGITransport
from app.main import app
from app.core.deps import get_db

@pytest.fixture
async def client(db_session):
    # zamijeni pravu bazu test sessionom
    app.dependency_overrides[get_db] = lambda: db_session
    transport = ASGITransport(app=app)
    async with AsyncClient(
        transport=transport, base_url="http://test"
    ) as c:
        yield c
    app.dependency_overrides.clear()
\end{lstlisting}

\vspace{0.4em}
\texttt{conftest.py} je poseban file - fixture iz njega su dostupne svim testovima bez importa. \texttt{dependency\_overrides} je FastAPI mehanizam: sve što je \texttt{Depends(get\_db)} u testu dobiva našu test sesiju umjesto prave baze.
\end{frame}

% ------------------------------------------------------------
\begin{frame}[fragile,shrink=5]{Izolacija testova: transakcija + rollback}
\begin{lstlisting}[language=python]
@pytest.fixture
async def db_session():
    async with engine.connect() as conn:
        trans = await conn.begin()
        session = AsyncSession(bind=conn)
        yield session          # test radi u transakciji
        await session.close()
        await trans.rollback() # sve promjene nestaju
\end{lstlisting}

\vspace{0.4em}
Svaki test dobiva svježu transakciju koja se na kraju rollbacka. Posljedice:
\begin{itemize}
  \item testovi ne vide podatke jedni od drugih
  \item redoslijed izvršavanja ne mijenja rezultat
  \item nema \enquote{očisti bazu} koda na kraju svakog testa
\end{itemize}

\vspace{0.4em}
Test koji ovisi o podacima prethodnog testa je tempirana bomba - radi dok netko ne pokrene samo jedan test.
\end{frame}

% ------------------------------------------------------------
\begin{frame}[fragile,shrink=5]{Fixture za tokene: testiranje kao admin i kao klub}
\begin{lstlisting}[language=python]
@pytest.fixture
async def admin_token(client):
    resp = await client.post("/auth/login", json={
        "username": "admin", "password": "admin123"})
    return resp.json()["access_token"]

@pytest.fixture
async def club_token(client):
    resp = await client.post("/auth/login", json={
        "username": "behemot", "password": "klub123"})
    return resp.json()["access_token"]

async def test_klub_ne_moze_kreirati_klub(client, club_token):
    resp = await client.post("/clubs",
        json={"name": "Novi klub"},
        headers={"Authorization": f"Bearer {club_token}"})
    assert resp.status_code == 403
\end{lstlisting}

\vspace{0.4em}
Za svaki endpoint pišemo barem jedan negativni test: pogrešna rola, tuđi resurs ili istekli rok.
\end{frame}

% ------------------------------------------------------------
\begin{frame}[fragile,shrink=5]{freezegun: vremenska pravila (rekapitulacija P6)}
\begin{lstlisting}[language=python]
from freezegun import freeze_time

@freeze_time("2027-03-05")  # izmedu prelim i final roka
async def test_nova_prijava_nakon_prelim_roka(client, club_token):
    resp = await client.post(
        "/competitions/1/registrations",
        json={"lifter_id": 1, "category": "93"},
        headers={"Authorization": f"Bearer {club_token}"})
    assert resp.status_code == 400
    assert resp.json()["code"] == "deadline_passed"

@freeze_time("2027-03-05")
async def test_odjava_izmedu_rokova_prolazi(client, club_token):
    resp = await client.post(
        "/competitions/1/registrations/1/withdraw",
        headers={"Authorization": f"Bearer {club_token}"})
    assert resp.status_code == 200
\end{lstlisting}

\vspace{0.4em}
\texttt{freeze\_time} zamrzava \texttt{datetime.now()} na zadani trenutak. Bez toga bi testovi faza prolazili ili padali ovisno o stvarnom datumu pokretanja.
\end{frame}

% ------------------------------------------------------------
\begin{frame}[fragile,shrink=5]{Coverage: koliko je koda pokriveno}
\begin{lstlisting}[language=bash]
uv pip install pytest-cov
pytest --cov=app --cov-report=term-missing

# Name                        Stmts   Miss  Cover   Missing
# app/services/registration.py  45      3    93%   88-90
# app/routers/clubs.py           30      0   100%
\end{lstlisting}

\vspace{0.4em}
\texttt{Missing} stupac pokazuje linije koje nijedan test ne izvršava - često su to error grane i edge caseovi, upravo ono što se najviše isplati testirati.

\vspace{0.6em}
Coverage mjeri \textbf{izvršenost}, ne ispravnost: linija može biti pokrivena testom koji ništa ne provjerava. 100\% coverage nije cilj - cilj je da su \textbf{poslovna pravila i error grane} pokriveni smislenim assertima. Za ovaj projekt: services i pravila blizu 100\%, boilerplate može niže.
\end{frame}

% ------------------------------------------------------------
\begin{frame}{Backend: što mora biti pokriveno prije obrane}
\begin{tabular}{@{}lp{7.2cm}@{}}
  \toprule
  \textbf{Područje} & \textbf{Minimum testova} \\
  \midrule
  Autentikacija & uspješan login, kriva lozinka, istekli token, refresh \\
  Autorizacija & svaka rola na svakom endpointu - tko smije, tko dobiva 403 \\
  Ownership & klub A ne vidi i ne mijenja resurse kluba B \\
  Poslovna pravila & svaka faza $\times$ svaka operacija iz matrice \\
  Constrainti & dupla prijava vraća 409, nevažeći input 422 \\
  \bottomrule
\end{tabular}

\vspace{0.6em}
Na obrani pitam: \enquote{Što točno testira ovaj test i što bi se moglo pokvariti da ga nema?} Test bez odgovora na to pitanje je balast.
\end{frame}

% ------------------------------------------------------------
\begin{frame}{Zašto testirati frontend}
Backend ima pytest - frontend mora imati ekvivalent. Bez testova:

\vspace{0.4em}
\begin{itemize}
  \item Bug u \texttt{parseValidationErrors} otkrije se na obrani, ne lokalno
  \item Refactoring Pinia storea razbije nešto u formi - ne znate dok ne kliknete
  \item CI pipeline može pustiti broken kod na produkciju
\end{itemize}

\vspace{0.6em}
\begin{tabular}{@{}lp{5cm}p{4cm}@{}}
  \toprule
  \textbf{Tip} & \textbf{Što testira} & \textbf{Alat} \\
  \midrule
  Unit & Izolirana funkcija ili store & Vitest \\
  Komponenta & Renderiranje s propsima & Vitest + test-utils \\
  E2E & Cijeli flow u browseru & Playwright \\
  \bottomrule
\end{tabular}
\end{frame}

% ------------------------------------------------------------
\begin{frame}[fragile,shrink=5]{Vitest: setup}
\begin{lstlisting}[language=bash]
# Instalacija
npm install -D vitest @vue/test-utils jsdom

# package.json - dodaj scripte
"test": "vitest",
"test:ui": "vitest --ui"
\end{lstlisting}

\begin{lstlisting}[language=typescript]
// vite.config.ts - dodaj test konfiguraciju
export default defineConfig({
  plugins: [vue()],
  test: {
    environment: 'jsdom',  // simulira browser DOM
    globals: true,         // describe/it/expect bez importa
  },
})
\end{lstlisting}

\vspace{0.4em}
Vitest koristi istu Vite konfiguraciju - nema zasebnog config filea. API je Jest-kompatibilan: \texttt{describe}, \texttt{it}, \texttt{expect}, \texttt{vi.mock}.
\end{frame}

% ------------------------------------------------------------
\begin{frame}{Što testirati unit testovima}
\begin{tabular}{@{}lp{5cm}p{4.5cm}@{}}
  \toprule
  \textbf{Isplati se testirati} & \textbf{Primjer} & \textbf{Zašto} \\
  \midrule
  Utility funkcije & \texttt{parseValidationErrors} & Čista logika, brzi testovi \\
  Pinia store actions & \texttt{fetchLifters} & Logika s mock API-jem \\
  Computed getteri & Filtrirani natjecatelji & Derivirana logika \\
  \bottomrule
\end{tabular}

\vspace{0.6em}
\textbf{Ne testirati:}
\begin{itemize}
  \item Vue framework sam (reaktivnost, direktive) - Vueov posao
  \item Third-party biblioteke (Axios, Pinia) - već imaju testove
  \item Čisti HTML prikaz bez logike - mala vrijednost, veliki maintenance
\end{itemize}
\end{frame}

% ------------------------------------------------------------
\begin{frame}[fragile,shrink=5]{Unit test: parseValidationErrors}
\begin{lstlisting}[language=typescript]
// src/utils/errors.test.ts
import { parseValidationErrors } from './errors'

describe('parseValidationErrors', () => {
  it('mapira FastAPI 422 detail na greske po polju', () => {
    const data = {
      detail: [
        { loc: ['body', 'first_name'],
          msg: 'field required', type: 'missing' },
        { loc: ['body', 'birth_date'],
          msg: 'invalid date', type: 'value_error' },
      ]
    }
    expect(parseValidationErrors(data)).toEqual({
      first_name: 'field required',
      birth_date: 'invalid date',
    })
  })

  it('vraca prazan objekt za ne-422 odgovor', () => {
    expect(parseValidationErrors({})).toEqual({})
  })
})
\end{lstlisting}
\end{frame}

% ------------------------------------------------------------
\begin{frame}[fragile,shrink=5]{Unit test: Pinia store}
\begin{lstlisting}[language=typescript]
// src/stores/lifters.test.ts
import { setActivePinia, createPinia } from 'pinia'
import { useLifterStore } from './lifters'
import api from '@/api/client'

vi.mock('@/api/client')  // zamijeni Axios s mockom

beforeEach(() => setActivePinia(createPinia()))

it('fetchLifters popunjava state uspjesno', async () => {
  const mockLifter = { id: 1, first_name: 'Ivan' }
  vi.mocked(api.get).mockResolvedValue(
    { data: [mockLifter] }
  )

  const store = useLifterStore()
  await store.fetchLifters(1)

  expect(store.lifters).toHaveLength(1)
  expect(store.lifters[0].first_name).toBe('Ivan')
  expect(store.loading).toBe(false)
  expect(store.error).toBeNull()
})
\end{lstlisting}
\end{frame}

% ------------------------------------------------------------
\begin{frame}[fragile,shrink=5]{Component test: FormField}
\begin{lstlisting}[language=typescript]
// src/components/FormField.test.ts
import { mount } from '@vue/test-utils'
import FormField from './FormField.vue'

it('prikazuje label i error poruku', () => {
  const wrapper = mount(FormField, {
    props: { label: 'Ime', error: 'Obavezno polje',
             modelValue: '' }
  })
  expect(wrapper.text()).toContain('Ime')
  expect(wrapper.find('.error').text()).toBe('Obavezno polje')
})

it('emitira update:modelValue pri unosu', async () => {
  const wrapper = mount(FormField, {
    props: { label: 'Ime', modelValue: '' }
  })
  await wrapper.find('input').setValue('Marko')
  expect(wrapper.emitted('update:modelValue')![0])
    .toEqual(['Marko'])
})
\end{lstlisting}

\vspace{0.4em}
\texttt{mount} renderira komponentu u jsdom-u. Testiramo \textbf{ugovor} komponente: propsi ulaze, eventi izlaze - ne unutarnju implementaciju.
\end{frame}

% ------------------------------------------------------------
\begin{frame}[fragile,shrink=5]{E2E testovi: Playwright}
Playwright kontrolira pravi browser - testira što korisnik zapravo vidi i radi:

\vspace{0.4em}
\begin{itemize}
  \item Pokreće Chromium, Firefox ili WebKit
  \item Navigira, klikće, ispunjava forme - kao pravi korisnik
  \item Provjera: URL, tekst na stranici, vidljivost elemenata
  \item Automatski čeka da elementi budu dostupni (nema sleep/timeout)
\end{itemize}

\vspace{0.6em}
\begin{lstlisting}[language=bash]
npm install -D @playwright/test
npx playwright install   # preuzme browsere
\end{lstlisting}

\vspace{0.4em}
U komponentama dodajte \texttt{data-testid} atribute za stabilne selectore: \texttt{data-testid="username"} ne puca kad promijenite CSS klasu ili tekst labela.
\end{frame}

% ------------------------------------------------------------
\begin{frame}[fragile,shrink=5]{E2E test: login flow}
\begin{lstlisting}[language=typescript]
// tests/e2e/login.spec.ts
import { test, expect } from '@playwright/test'

test('prijava preusmjerava na dashboard', async ({ page }) => {
  await page.goto('/login')

  await page.getByTestId('username').fill('behemot')
  await page.getByTestId('password').fill('klub123')
  await page.getByTestId('submit').click()

  await expect(page).toHaveURL('/dashboard')
  await expect(
    page.getByText('Dobro dosli')
  ).toBeVisible()
})

test('pogresna lozinka prikazuje gresku', async ({ page }) => {
  await page.goto('/login')
  await page.getByTestId('username').fill('behemot')
  await page.getByTestId('password').fill('krivo')
  await page.getByTestId('submit').click()
  await expect(page.getByText('Pogresni podaci')).toBeVisible()
})
\end{lstlisting}
\end{frame}

% ------------------------------------------------------------
\begin{frame}[fragile,shrink=5]{E2E test: pravilo faze u UI-ju}
Seed sadrži natjecanje kojem je final rok prošao. E2E test provjerava da UI poštuje fazu:

\begin{lstlisting}[language=typescript]
// tests/e2e/phases.spec.ts
import { test, expect } from '@playwright/test'
import { loginAsClub } from './helpers'

test('nakon final roka gumb odjave je disabled',
    async ({ page }) => {
  await loginAsClub(page)
  await page.goto('/competitions/2/registrations')

  await expect(
    page.getByTestId('withdraw-btn')
  ).toBeDisabled()
})
\end{lstlisting}

\vspace{0.4em}
Ovaj test spaja sve: backend vraća faze, frontend computed ih pretvara u \texttt{:disabled}, a test dokazuje da korisnik nakon roka stvarno ne može kliknuti. Unit test to ne može pokriti - zato E2E postoji.
\end{frame}

% ------------------------------------------------------------
\begin{frame}[fragile,shrink=5]{Playwright config: automatski dev server}
\begin{lstlisting}[language=typescript]
// playwright.config.ts
import { defineConfig } from '@playwright/test'

export default defineConfig({
  testDir: 'tests/e2e',
  use: { baseURL: 'http://localhost:5173' },
  webServer: {
    command: 'npm run dev',
    port: 5173,
    reuseExistingServer: true,
  },
})
\end{lstlisting}

\vspace{0.4em}
\texttt{webServer} pokreće Vite dev server prije testova i ugasi ga nakon njih. Backend i baza moraju biti pokrenuti zasebno - E2E testira protiv pravog API-ja sa seed podacima.

\vspace{0.4em}
Zato E2E testove pokrećemo protiv \textbf{poznatog stanja}: \texttt{python -m app.seed} prije test runa vraća bazu u početno stanje.
\end{frame}

% ------------------------------------------------------------
\begin{frame}{Testovi u CI pipelineu}
\begin{tabular}{@{}lp{4cm}p{5cm}@{}}
  \toprule
  \textbf{Test} & \textbf{Kad se pokreće} & \textbf{Zašto} \\
  \midrule
  pytest & Svaki push & Poslovna pravila backenda \\
  Vitest & Svaki push & Brz (sekunde), hvatamo regresije rano \\
  Playwright & Pull request & Sporiji (minute), treba running server \\
  \bottomrule
\end{tabular}

\vspace{0.6em}
\begin{itemize}
  \item \texttt{pytest} i \texttt{npm run test -- --run} idu u svaki CI run
  \item Ako test padne - commit dobije crveni X, merge je blokiran
  \item Cijeli pipeline (lint + test + build) slažemo na Predavanju 12
\end{itemize}

\vspace{0.4em}
Za studentski projekt obavezni su pytest i Vitest. Playwright je bonus koji pokazuje zrelost projekta na obrani.
\end{frame}

% ------------------------------------------------------------
\begin{frame}{Commitovi za Predavanje 11}
\begin{itemize}
  \item \texttt{test(api): add conftest with async client and db fixtures}
  \item \texttt{test(api): add coverage config and fill gaps in rule tests}
  \item \texttt{test(web): add vitest setup with util and store unit tests}
  \item \texttt{test(web): add component tests for form field}
  \item \texttt{test(e2e): add playwright login and phase rule tests}
\end{itemize}

\vspace{0.4em}
\textbf{Na satu:} pokrećemo testove, rušimo jedno poslovno pravilo namjerno i gledamo koji test pada i zašto.
\end{frame}

% ------------------------------------------------------------
\begin{frame}{Što ide sljedeće (Predavanje 12)}
\begin{itemize}
  \item Production konfiguracija i environment varijable
  \item CORS i deploy backenda na Railway
  \item Deploy frontenda i SPA fallback
  \item CI/CD pipeline koji pokreće današnje testove
\end{itemize}

\vspace{0.8em}
\begin{center}
\Large Pitanja?
\end{center}
\end{frame}

\end{document}
